\section{Descrição do Processo}\label{sec:processo}

\subsection{Configuração Geral da UTAA}

O esquema de Tratamento de Águas Ácidas considerada neste trabalho é o esquema composto por duas torres de destilação operando em série.

Nesse esquema de tratamento, a água ácida proveniente de diversas unidades da refinaria (hidrotratamento, FCC, coqueamento, etc.) é coletada e alimentada na primeira torre, a torre esgotadora de H\textsubscript{2}S. Esta coluna opera sob pressão moderada e utiliza vapor de esgotamento e um refervedor para promover a evaporação preferencial do H\textsubscript{2}S, que é mais volátil que a NH\textsubscript{3} nas condições operacionais típicas. O gás ácido de topo, rico em H\textsubscript{2}S e contendo alguma NH\textsubscript{3}, é enviado para a URE. A água de fundo, ainda contendo NH\textsubscript{3} e traços de H\textsubscript{2}S, é bombeada para a segunda torre, onde a NH\textsubscript{3} é removida.


\subsection{Desafios Operacionais}

A operação da torre esgotadora apresenta os seguintes desafios:

\begin{enumerate}
  \item \textbf{Maximização da recuperação de H\textsubscript{2}S}: É desejável maximizar a concentração de H\textsubscript{2}S no gás ácido para otimizar a operação da URE.
  \item \textbf{Minimização de NH\textsubscript{3} no gás ácido}: A presença de NH\textsubscript{3} no gás ácido leva à formação de sais de amônio (principalmente NH\textsubscript{4}HS), que causam corrosão, entupimentos e redução de eficiência na URE.
  \item \textbf{Prevenção de fervura da coluna}: Em condições extremas de carga térmica ou composição, pode ocorrer o fenômeno de "fervura", caracterizado por um aumento abrupto da taxa de vaporização que leva a arraste excessivo de líquido e NH\textsubscript{3} para o topo.
  \item \textbf{Eficiência energética}: A carga térmica do refervedor representa um custo operacional significativo, sendo desejável minimizá-la mantendo os objetivos de separação.
\end{enumerate}


\subsection{Simulação Dinâmica de Colunas de Destilação}

A simulação dinâmica de colunas de destilação tem sido amplamente estudada na literatura, com ênfase em aspectos como convergência numérica e inicialização de modelos. Trabalhos como o de Secchi et al.\ (1993) \cite{secchi1993} apresentam modelagens matemáticas rigorosas para processos de destilação, destacando a necessidade de métodos numéricos eficientes para resolver as equações diferenciais e algébricas que descrevem o comportamento dinâmico das colunas.

A inicialização de modelos dinâmicos a partir de soluções estacionárias é uma prática consolidada para aumentar a robustez e acelerar a convergência. Conforme discutido em modelos simplificados para colunas binárias \cite{vieira2001}, perfis estacionários de composição, temperatura e pressão obtidos por simuladores estáticos podem ser transferidos para o modelo dinâmico, reduzindo discrepâncias iniciais e mitigando o risco de divergência do integrador. Essa estratégia é especialmente recomendada para modelos de equilíbrio de estágio, em que a solução estacionária fornece um ponto de partida fisicamente consistente para a simulação transitória.




\subsection{Modelagem Dinâmica no Aspen Plus Dynamics V14}\label{sec:modelo_dinamico}

\subsubsection{Preparação do Modelo Estático}

Antes da exportação para o ambiente dinâmico, o modelo estático foi preparado conforme os seguintes passos:

\begin{enumerate}
  \item Verificação de convergência robusta do modelo estático
  \item Dimensionamento de equipamentos (diâmetro de coluna, volumes de pratos, holdups)
  \item Especificação de controladores básicos de inventário (nível, pressão)
  \item Definição de válvulas de controle e suas características
  \item Configuração de instrumentação (transmissores, sensores)
\end{enumerate}

\subsubsection{Exportação e Configuração em Aspen Dynamics}

O modelo foi exportado do Aspen Plus para o Aspen Plus Dynamics utilizando a funcionalidade de exportação de modelos em regime permanente (\emph{Pressure-Driven} mode). As principais configurações incluem:

leftmargin=4em, itemsep=0pt
  \item Modo de simulação: modo orientado por pressão (\emph{pressure-driven})
  \item Passo de integração: [valor] segundos
  \item Método de integração: [especificar]
  \item Configuração de controladores PID para inventários
  \item Especificação de dinâmicas de válvulas
\end{itemize}

\subsubsection{Condições Iniciais e Convergência Dinâmica}

A inicialização do modelo dinâmico utilizou o estado estacionário convergido do modelo estático. As seguintes verificações foram realizadas:

\begin{enumerate}
  \item Balanços de massa e energia em cada estágio
  \item Estabilidade de holdups e taxas de fluxo
  \item Resposta de controladores de inventário
  \item Tempo de estabelecimento após pequenas perturbações
\end{enumerate}



\subsection{Métricas de Desempenho}

Para avaliar o desempenho operacional da torre em cada cenário, as seguintes métricas serão calculadas:

leftmargin=4em, itemsep=0pt
  \item \textbf{Recuperação de H\textsubscript{2}S}: fração mássica de H\textsubscript{2}S da alimentação que sai no gás ácido de topo
  \item \textbf{Concentração de NH\textsubscript{3} no gás ácido}: ppm (mol) de NH\textsubscript{3} no topo
  \item \textbf{Consumo energético}: carga térmica do refervedor por unidade de massa de água ácida tratada
  \item \textbf{Margem até fervura}: distância operacional até condições críticas
  \item \textbf{Qualidade da água tratada}: concentração residual de H\textsubscript{2}S e NH\textsubscript{3} no fundo
\end{itemize}









\subsection{Modelo Dinâmico: Resposta Transitória}

[Apresentar resultados de simulações dinâmicas para perturbações típicas]

% \begin{figure}[H]
%   \centering
%   % \includegraphics[width=0.8\linewidth]{resposta_dinamica.png}
%   \caption{Resposta dinâmica a perturbação em degrau na composição de alimentação.}
%   \label{fig:resposta_dinamica}
% \end{figure}


\section{Conclusões}\label{sec:conclusoes}

% TODO: Sintetizar principais conclusões do trabalho

Este trabalho apresentou o desenvolvimento e análise de modelos de simulação estáticos e dinâmicos de uma Unidade de Tratamento de Águas Ácidas (UTAA), com foco particular na torre esgotadora de H\textsubscript{2}S. Os modelos foram implementados em Aspen Plus V14 e Aspen Plus Dynamics V14, utilizando o pacote termodinâmico GPSWAT para representação adequada dos equilíbrios físico-químicos envolvendo eletrólitos em solução aquosa.

[Discutir principais conclusões sobre:]
leftmargin=4em, itemsep=0pt
  \item Capacidade dos modelos de reproduzir o comportamento esperado do processo
  \item Trade-offs identificados entre recuperação de H\textsubscript{2}S, minimização de NH\textsubscript{3} e eficiência energética
  \item Caracterização do ponto operacional crítico (limite de fervura)
  \item Sensibilidade do sistema a variações de composição e condições operacionais
  \item Implicações para estratégias de operação e controle
\end{itemize}

Os modelos desenvolvidos fornecem uma base sólida para otimização operacional da UTAA e para o desenvolvimento de estratégias de controle avançado, conforme discutido na seção de trabalhos futuros.

\section{Trabalhos Futuros}\label{sec:futuros}

Este trabalho estabelece as fundações para desenvolvimentos subsequentes em controle avançado da UTAA. As seguintes direções são propostas:

\subsection{Geração de Dados para Treinamento de Redes Neurais}

O modelo dinâmico desenvolvido em Aspen Plus Dynamics será utilizado para gerar extensos conjuntos de dados de treinamento mediante simulação de trajetórias operacionais diversificadas. Especificamente:

leftmargin=4em, itemsep=0pt
  \item Simulação de múltiplos cenários com variações nas composições de alimentação, vazões e condições operacionais
  \item Coleta de séries temporais de variáveis de processo (temperaturas, composições, pressões, fluxos)
  \item Identificação de relações entrada-saída para mapeamento via redes neurais
  \item Geração de dados próximos ao limite de fervura para capturar não-linearidades críticas
\end{itemize}

\subsection{Desenvolvimento de Modelos de Redes Neurais}

Com base nos dados gerados, redes neurais artificiais (ANNs) serão treinadas para emular o comportamento dinâmico da torre esgotadora. Arquiteturas candidatas incluem:

leftmargin=4em, itemsep=0pt
  \item Redes neurais feedforward (FNNs) para mapeamento estático entrada-saída
  \item Redes neurais recorrentes (RNNs, LSTMs, GRUs) para captura de dinâmicas temporais
  \item Modelos híbridos combinando conhecimento fenomenológico com componentes baseados em dados
\end{itemize}

\subsection{Controle Preditivo Não Linear Baseado em Redes Neurais (NN NMPC)}

O objetivo final é desenvolver um controlador preditivo não linear (NMPC) utilizando o modelo de rede neural como preditor interno. A estrutura proposta envolve:

\begin{enumerate}
  \item \textbf{Modelo de predição}: Rede neural treinada substituindo o modelo fenomenológico rigoroso no otimizador do NMPC, reduzindo drasticamente o custo computacional
  \item \textbf{Função objetivo}: Minimização de um funcional custo considerando rastreamento de referências (recuperação de H\textsubscript{2}S), penalização de NH\textsubscript{3} no topo, esforço de controle e eficiência energética
  \item \textbf{Restrições}: Limites operacionais em variáveis manipuladas (carga térmica, vazão de vapor) e variáveis de processo (temperatura, pressão, concentrações)
  \item \textbf{Estimação de estados}: Filtros estendidos de Kalman (EKF) ou estimadores por horizonte móvel (MHE) para estados não medidos
  \item \textbf{Operação offset-free}: Incorporação de estados integradores para rejeição de distúrbios e desalinhamentos modelo-planta
\end{enumerate}

\subsection{Validação e Implementação}

O controlador NN NMPC será validado inicialmente em simulações fechadas (closed-loop) utilizando o modelo rigoroso do Aspen Plus Dynamics como "planta virtual". Serão avaliados:

leftmargin=4em, itemsep=0pt
  \item Desempenho de rastreamento de setpoints
  \item Rejeição de distúrbios em composição e vazão de alimentação
  \item Capacidade de operação próxima ao limite de fervura sem violação de restrições
  \item Robustez a incertezas de modelo e ruídos de medição
  \item Comparação com estratégias de controle convencionais (PID, MPC linear)
\end{itemize}

Caso os resultados sejam promissores, estudos de viabilidade para implementação em planta piloto ou industrial poderão ser conduzidos.




\subsubsection{Síntese e Implicações Operacionais}

A análise de sensibilidade realizada fornece subsídios valiosos para a compreensão do comportamento da torre esgotadora de H\textsubscript{2}S e para o desenvolvimento de estratégias de operação e controle. Os principais achados são sintetizados a seguir:

\begin{itemize}[leftmargin=4em, itemsep=0pt]
  \item \textbf{Trade-off fundamental}: Existe uma região operacional ótima onde se atinge recuperação completa de H\textsubscript{2}S com mínima perda de NH\textsubscript{3}. Operar abaixo dessa região resulta em perda de H\textsubscript{2}S (redução de eficiência da URE), enquanto operar acima resulta em arraste excessivo de NH\textsubscript{3} (risco de formação de sais de amônio na URE).

  \item \textbf{Concentração de contaminantes}: Alimentações com alta concentração de NH\textsubscript{3} apresentam maior propensão ao arraste de amônia em cargas térmicas elevadas. A operação com alimentações diluídas oferece maior flexibilidade operacional e menor risco de perda de NH\textsubscript{3}.

  \item \textbf{Temperatura de alimentação}: O pré-aquecimento da alimentação reduz significativamente o consumo energético do refervedor (até 10-15\%), porém aumenta o risco de arraste de NH\textsubscript{3} se não for acompanhado de ajuste adequado na carga térmica. A integração energética deve ser avaliada em conjunto com os requisitos de qualidade do gás ácido.

  \item \textbf{Vazão de alimentação}: Variações na vazão de alimentação demandam ajustes proporcionais na carga térmica do refervedor. A relação carga térmica/vazão (energia específica) é uma variável chave para evitar arraste excessivo de NH\textsubscript{3} em condições de baixa vazão ou perda de H\textsubscript{2}S em condições de alta vazão.

  \item \textbf{Temperatura de fundo como indicador}: A temperatura de fundo da torre mostrou-se um indicador confiável da proximidade ao ponto de arraste crítico de NH\textsubscript{3}. Temperaturas de fundo superiores a 160 °C devem ser evitadas para minimizar perda de amônia, sugerindo que o controle da temperatura de fundo é uma estratégia viável.
\end{itemize}

Esses resultados estabelecem a base para o desenvolvimento de estratégias de controle preditivo não linear (NMPC) utilizando redes neurais, conforme objetivos futuros do projeto. Os modelos dinâmicos desenvolvidos em Aspen Plus Dynamics V14 permitirão gerar datasets para treinamento de redes neurais que capturem a dinâmica não linear do sistema, viabilizando controle avançado em tempo real considerando múltiplas variáveis manipuladas (carga térmica do refervedor, temperatura de alimentação) e múltiplos objetivos (maximização de recuperação de H\textsubscript{2}S, minimização de perda de NH\textsubscript{3}, minimização de consumo energético).



Conclusao



Os resultados obtidos evidenciam também uma implicação importante para futuras estratégias de controle da torre esgotadora. A análise das curvas revela a existência de dois regimes operacionais distintos, caracterizados por variações significativas no ganho estático entre a carga térmica do refervedor e a temperatura de topo. Essa mudança estrutural implica que a dinâmica da coluna não permanece aproximadamente linear em toda a faixa de operação. Como consequência, uma estratégia de controle PID implementada de forma direta e sintonizada apenas para o regime de alta sensibilidade (associado ao aumento rápido da recuperação de H\textsubscript{2}S) pode apresentar desempenho degradado ou mesmo tornar-se instável quando a operação se desloca para a região de baixa sensibilidade, observada após a recuperação completa de H\textsubscript{2}S. Esse comportamento reforça a necessidade de abordagens de controle mais robustas ou adaptativas em aplicações industriais que envolvem variações significativas nas condições de contaminação ou na carga térmica aplicada.
